\subsection{Within-District Partisan Homogeneity}

Table~\ref{tab:psd-homogeneity} reports mean within-district partisan
standard deviation ($\sigma_D$) for North Carolina and Texas across four
$\alpha$ levels, compared to the enacted plan and the neutral bisect baseline.

\begin{table}[ht]
\centering
\caption{Mean within-district partisan standard deviation ($\sigma_D$) by
  $\alpha$ level.\singrun}
\label{tab:psd-homogeneity}
\begin{tabular}{llr}
\toprule
State & Map & Mean $\sigma_D$ (pp) \\
\midrule
\multirow{6}{*}{NC ($k=14$)}
& Neutral bisect ($\alpha = 0$) & 14.8 \\
& PSD $\alpha=5$     & 12.1 \\
& PSD $\alpha=20$    &  9.3 \\
& PSD $\alpha=50$    &  7.8 \\
& Enacted 2020 plan  & 11.4 \\
\midrule
\multirow{6}{*}{TX ($k=38$)}
& Neutral bisect ($\alpha = 0$) & 16.2 \\
& PSD $\alpha=5$     & 14.3 \\
& PSD $\alpha=20$    & 11.1 \\
& PSD $\alpha=50$    &  9.2 \\
& Enacted 2020 plan  & 12.8 \\
\bottomrule
\end{tabular}
\end{table}

PSD weighting at $\alpha = 50$ reduces $\sigma_D$ from 14.8 to 7.8 pp in
North Carolina and from 16.2 to 9.2 pp in Texas---reductions of 47\% and
43\% respectively.\singrun  Crucially, the enacted plans (with $\sigma_D =
11.4$ pp and $12.8$ pp) have lower $\sigma_D$ than the neutral baseline but
substantially higher than the PSD maps at $\alpha \geq 20$.  This shows
that: (a) the enacted plans do contain deliberate partisan clustering beyond
what geography alone produces, and (b) algorithmically generated PSDs can
achieve substantially more homogeneity than even aggressive enacted
gerrymanders.

\subsection{Safe Seat Distribution}

Table~\ref{tab:psd-safeseats} reports safe seat counts at the $\alpha = 20$
level, where the homogeneity benefit stabilises.

\begin{table}[ht]
\centering
\caption{Safe seat counts by party ($\alpha=20$, PSD vs.\ enacted).\singrun}
\label{tab:psd-safeseats}
\begin{tabular}{llrrrr}
\toprule
State & Map & Safe D ($D\%>60$) & Safe R ($D\%<40$) & Competitive & EG \\
\midrule
NC & Neutral bisect & 4 &  5 & 5 & $+0.03$ \\
NC & PSD $\alpha=20$ & 5 & 6 & 3 & $-0.01$ \\
NC & Enacted 2020   & 3 & 8 & 3 & $+0.14$ \\
\midrule
TX & Neutral bisect & 12 & 18 &  8 & $+0.06$ \\
TX & PSD $\alpha=20$ & 15 & 19 &  4 & $+0.02$ \\
TX & Enacted 2020   &  9 & 23 &  6 & $+0.18$ \\
\bottomrule
\end{tabular}
\end{table}

The PSD maps distribute safe seats more symmetrically between the parties
than the enacted plans: in North Carolina, the PSD at $\alpha=20$ produces
5 safe Democratic and 6 safe Republican districts (ratio 5:6, close to the
statewide vote share ratio), while the enacted plan produces 3:8 (strongly
asymmetric).\singrun  In Texas, the PSD produces 15:19 (close to the 40/60
statewide partisan split) while the enacted plan produces 9:23.

The Efficiency Gap confirms this pattern.  Neutral bisect maps have a small
positive EG ($+0.03$ to $+0.06$), reflecting the Rodden gap (geographic
sorting slightly advantages Republicans).  PSD maps have an EG near zero
($-0.01$ to $+0.02$), because the algorithm distributes similar-leaning tracts
symmetrically regardless of party.  The enacted plans have EGs of $+0.14$ and
$+0.18$, well above the $\pm 0.08$ threshold that \citet{stephanopoulos2015partisan}
proposed as a gerrymandering indicator.

\subsection{Compactness Cost}

Table~\ref{tab:psd-compactness} reports mean Polsby-Popper as $\alpha$
increases.

\begin{table}[ht]
\centering
\caption{Compactness vs.\ partisan similarity: mean Polsby-Popper.\singrun}
\label{tab:psd-compactness}
\begin{tabular}{llr}
\toprule
State & Map & Mean PP \\
\midrule
NC & Neutral bisect ($\alpha=0$) & 0.171 \\
NC & PSD $\alpha=5$    & 0.164 \\
NC & PSD $\alpha=20$   & 0.148 \\
NC & PSD $\alpha=50$   & 0.121 \\
NC & Enacted 2020      & 0.143 \\
\midrule
TX & Neutral bisect ($\alpha=0$) & 0.198 \\
TX & PSD $\alpha=5$    & 0.188 \\
TX & PSD $\alpha=20$   & 0.171 \\
TX & PSD $\alpha=50$   & 0.139 \\
TX & Enacted 2020      & 0.162 \\
\bottomrule
\end{tabular}
\end{table}

Partisan similarity weighting degrades compactness.  At $\alpha = 20$, mean
PP falls by 13\% in NC and 14\% in TX relative to the neutral baseline.\singrun
At $\alpha = 50$, the fall is 29\% and 30\% respectively.  This is expected:
when the algorithm clusters partisan-similar tracts, it may combine
geographically non-adjacent (but politically similar) tracts, producing
elongated or irregularly shaped districts.

Notably, the enacted plans have PP values ($0.143$ and $0.162$) close to the
PSD at $\alpha = 20$ ($0.148$ and $0.171$).  This suggests that enacted
partisan gerrymanders operate at an implicit partisan-similarity weighting
roughly equivalent to $\alpha \approx 15$--$20$---moderate partisan clustering
with a meaningful compactness cost but not the extreme homogeneity achievable
at higher $\alpha$.

\subsection{Geographic Distribution of Safe Seats}

PSD safe seats are spatially contiguous clusters of similar partisan precincts:
North Carolina's PSD Democratic safe districts correspond to the Research
Triangle (Durham-Chapel Hill-Raleigh) and Charlotte urban core, while
Republican safe districts cover the Piedmont and western mountain counties.
This spatial pattern mirrors geographic sorting closely---PSDs make explicit
what sorting implies implicitly.

The enacted gerrymandered map in NC splits the Research Triangle into multiple
districts, reducing the number of safe Democratic seats from 5 (PSD) to 3
(enacted), while consolidating Republican-leaning suburban precincts with
surrounding rural territory to produce 8 safe Republican districts.  The PSD
analysis quantifies this manipulation: the excess Republican safe seats
(8 enacted vs.\ 6 PSD) represent the incremental effect of intentional
boundary manipulation beyond geographic sorting.
